\chapter{Long-Haul Fibre}\label{ch:nw:long-haul-fibre}

In August 2010 a new fibre route between Chicago and New Jersey went into service: shorter than the routes before it, with lower-latency
equipment, and quoted at \qty{6.65}{\milli\second} one way between the Cermak building in Chicago and Nasdaq's data centre in Carteret,
some \qty{13.3}{\milli\second} round trip. Its cost was publicly estimated at about 300~million dollars. Before it, prices in New Jersey had
responded to Chicago within 7.25 to \qty{7.95}{\milli\second}; around its opening, the response times drop by 1 to
\qty{1.5}{\milli\second}. Within two years microwave
networks estimated at 4.2 to \qty{5.2}{\milli\second} had beaten it, and the cable's buyers had learnt what chapter~10's floors say: glass is
slower than air by 46\%, and no amount of digging changes that.

Fibre did not lose. It carries everything that is not the fastest few messages: market data by the gigabit, orders when the weather closes
the radio links (chapter~14), the transatlantic routes that no tower can cross, and every connection from a firm's cabinet to its carrier.
This chapter takes a long-haul route apart: the glass and what happens to light in it, the path and how straight it is, the equipment, the
subsea cables, the ways of buying fibre, and the new glass that is nearly as fast as air.

\section{Light in glass: attenuation, amplifiers and dispersion}

Standard single-mode fibre at \qty{1550}{\nano\metre} loses about \qty{0.16}{\decibel} per kilometre (Corning's datasheet for its
ultra-low-loss fibre gives 0.16 typical, 0.17 maximum), so a signal loses half its power every 19~kilometres and a
thousandth of it every 190. Over the \qty{1359}{\kilo\metre} of glass that the 2010 route must contain (below), the loss would be
\qty{217}{\decibel}: no receiver hears that. Long-haul fibre is therefore cut into spans, with an amplifier at the end of each.

\begin{definition}[Optical amplifier]\label{def:nw:long-haul-fibre:amp}
An \emph{optical amplifier}\index{optical amplifier} restores the power of the light in a fibre without converting it to an electrical
signal, typically by passing it through a few metres of fibre doped with a rare earth and pumped by a laser; placed at the end of each span,
it lets a signal cross thousands of kilometres of glass.
\end{definition}

The amplifier costs little time (the doped fibre is metres, the pass-through a fraction of a microsecond), which is why the route model below
charges \qty{0.1}{\micro\second} per span, an assumption, and why amplification is not the latency problem. Dispersion is.

\begin{definition}[Chromatic dispersion, dispersion compensation]\label{def:nw:long-haul-fibre:disp}
\emph{Chromatic dispersion}\index{chromatic dispersion} is the spreading of a light pulse in a fibre because its different wavelengths travel
at slightly different speeds, measured in picoseconds of spread per nanometre of spectral width per kilometre. \emph{Dispersion
compensation}\index{dispersion compensation} undoes it, either optically, with a length of fibre of opposite dispersion (dispersion-compensating
fibre) in each span, or electronically, by signal processing in a coherent receiver.
\end{definition}

The same datasheet bounds dispersion at \qty{18}{\pico\second} per nanometre per kilometre at \qty{1550}{\nano\metre}: over 1\,359~kilometres, a
spread of about \qty{24}{\nano\second} per nanometre of spectrum, far more than a bit at ten gigabits (\qty{0.1}{\nano\second}). Compensating it
optically adds glass to the path: a spool in every span, which the model sets at a fifth of the span's length. Compensating it electronically
adds processing time at the ends instead, and leaves the path alone.

\begin{definition}[Forward error correction]\label{def:nw:long-haul-fibre:fec}
\emph{Forward error correction}\index{forward error correction} (FEC) adds redundant bits to a transmitted signal so that the receiver can
correct errors without asking for a retransmission; its decoder must collect a block (or an interleaved group of blocks) before it can correct
it, which adds a delay at the receiving end.
\end{definition}

FEC lets a link run with a noisier signal, hence longer spans or higher speeds, at the price of a delay per terminal. Laughlin, Aguirre and
Grundfest note the same trade-off in microwave radios, whose latency is ``dominated by'' FEC buffering and interleaving. For low-latency
fibre the choices are therefore the ones the model exposes: short spans or long ones, optical or electronic \omterm{def:nw:long-haul-fibre:disp}{dispersion compensation}, heavy
FEC or light.

\begin{omfigure}
\begin{tikzpicture}[font=\footnotesize,b/.style={draw,rounded corners=2pt,minimum height=0.6cm,align=center,inner sep=3pt},>=Stealth]
\node[b,fill=omThm!15] (t1) at (0,0) {terminal:\\ transponder\\ and FEC};
\node[b,fill=omDef!15] (a1) at (4.6,0) {amplifier};
\node[b,fill=omDef!15] (a2) at (8.4,0) {amplifier};
\node[b,fill=omThm!15] (t2) at (12.4,0) {terminal:\\ FEC decode\\ and transponder};
\draw[thick] (t1) -- node[above=2pt]{span 1: 80~km} (a1);
\draw[thick] (a1) -- node[above]{span 2} (a2);
\draw[thick,dashed] (a2) -- node[above]{\dots\ span $N$} (t2);
\draw[omMeth,thick] (6.5,-0.95) circle (0.28) node[below=0.3cm,font=\scriptsize,omMeth]{optional DCF spool};
\draw[omMeth,->] (6.5,-0.67) -- (6.5,-0.05);
\node[font=\scriptsize,align=center] at (6.2,1.1) {latency = glass along the geodesic + extra path + DCF + amplifiers + terminals};
\end{tikzpicture}

\omcaption{A long-haul fibre route, schematically: terminals at each end (transponders and \omterm{def:nw:long-haul-fibre:fec}{forward error correction}), spans of glass with an
amplifier at the end of each, and, on routes that compensate dispersion optically, a spool of dispersion-compensating fibre (DCF) in each span.
\texttt{firm.fibreroute} charges each piece its time.}\label{fig:nw:long-haul-fibre:span}
\end{omfigure}

\section{Route straightness: the Chicago--New York story}

Chapter~10's floor for 350 Cermak to Carteret is \qty{5507}{\micro\second} in straight fibre; the 2010 route was quoted at
\qty{6650}{\micro\second}. The difference, \qty{1143}{\micro\second}, is path and equipment, and inverting the published figure separates them
once an equipment time is assumed. With two terminals of \qty{10}{\micro\second} and an amplifier every 80~kilometres at
\qty{0.1}{\micro\second} (the model's assumptions, \qty{21.7}{\micro\second} in all), the route must contain 1\,359~kilometres of glass: 230
kilometres more than the geodesic, a path factor of 1.20. The later figure of \qty{6.55}{\milli\second} implies 1\,339~kilometres, 20 fewer,
or faster equipment.

\begin{proposition}[Inverting a published latency]\label{prop:nw:long-haul-fibre:invert}
If a fibre route of \omterm{def:nw:the-north-american-map:floor}{group index} $n_g$ is published at $T$ one way and its equipment takes $E$, its glass is
\[
  L = \frac{(T - E)\,c_0}{n_g},
\]
and its path factor is $L/d$ for a geodesic $d$. An error $\delta E$ in the assumed equipment time moves $L$ by $\delta E\,c_0/n_g$, about
\qty{205}{\metre} per microsecond.
\end{proposition}

\begin{proof}
Light covers $L$ at $c_0/n_g$ in the time left once the equipment's $E$ is removed. Differentiating in $E$ gives the sensitivity.
\end{proof}

\omcode{code/firm/fibreroute/firm_fibreroute.py}{59}{75}{A route's latency by component, and the inversion of a published latency into the
glass the route must contain.}\label{lst:nw:long-haul-fibre:components}

\begin{omfigure}
\begin{tikzpicture}
\begin{axis}[width=11cm,height=5.6cm,xmin=0,xmax=200,ymin=1.16,ymax=1.215,xlabel={equipment time assumed (\si{\micro\second}, one way)},
  ylabel={implied path factor},tick label style={font=\footnotesize},label style={font=\small},grid=major,grid style={gray!25},
  table/col sep=comma,yticklabel style={/pgf/number format/fixed,/pgf/number format/precision=2}]
\addplot[omDef,thick,no markers] table[x=equipment_us,y=factor]{figdata/networks/13-long-haul-fibre/sensitivity.csv};
\end{axis}
\end{tikzpicture}

\omcaption{How much of the 2010 route's excess is path: the path factor implied by its quoted \qty{6.65}{\milli\second} as a function of the
equipment time assumed. Even \qty{200}{\micro\second} of equipment leaves a path 17\% longer than the geodesic. Data: \texttt{fig\_fibre.py},
\texttt{nw\_fibre.sensitivity()}.}\label{fig:nw:long-haul-fibre:sensitivity}
\end{omfigure}

\Cref{fig:nw:long-haul-fibre:sensitivity} shows why the conclusion is robust: whatever the equipment, most of the millisecond is path. A cable
must follow rights of way it can obtain; the 2010 route bought straightness at 300~million dollars and still ran 17 to 20\% longer than the
geodesic. The microwave networks that beat it had \omterm{def:nw:the-north-american-map:factor}{route factors} near 1.01 in air (chapter~10) and a medium 46\% faster.

\begin{table}[ht]
\centering\small
\begin{tabular}{lrrrrr}
\toprule
Route & Geodesic & Published & Fibre floor & Glass & Path \\
 & (km) & (ms, one way) & (ms) & (km) & factor \\
\midrule
Chicago--New Jersey, 2010 & 1\,129.2 & 6.650 & 5.507 & 1\,359 & 1.204 \\
Chicago--New Jersey, later & 1\,129.2 & 6.550 & 5.507 & 1\,339 & 1.185 \\
New York--London subsea, 2015 & 5\,551.0 & 29.475 & 27.070 & 6\,038 & 1.088 \\
\bottomrule
\end{tabular}
\caption{Published fibre routes inverted into glass and path factor (\cref{prop:nw:long-haul-fibre:invert}), with the model's equipment
assumptions (\qty{10}{\micro\second} per terminal, \qty{0.1}{\micro\second} per 80-km span): Laughlin, Aguirre and Grundfest's quoted figures
for the Chicago route, half of Hibernia Express's published round trip for the Atlantic. Data: \texttt{nw\_fibre.published()}.}\label{tab:nw:long-haul-fibre:published}
\end{table}

\section{Subsea cables}

\begin{definition}[Cable landing station]\label{def:nw:long-haul-fibre:cls}
A \emph{cable landing station}\index{cable landing station} is the building where a submarine cable comes ashore and connects to terrestrial
networks: it powers the cable's undersea repeaters, terminates its fibre pairs and hands their traffic to land routes.
\end{definition}

A subsea route is a terrestrial route at each end, a landing station on each coast, and the ocean between. Hibernia Express, the first new
transatlantic cable in over twelve years when it opened in 2015, published a round trip under \qty{58.95}{\milli\second} between Equinix's NY4
in Secaucus and LD4 in Slough, more than five milliseconds better than the existing cables, over a fibre pair of Corning's pure-silica-core
fibre. The geodesic between the two buildings is \qty{5551}{\kilo\metre}; half the round trip, \qty{29.475}{\milli\second}, implies about
6\,040~kilometres of glass, a path factor of 1.09. Across an ocean the geodesic is almost available, apart from the landings and the tails to the
data centres; on land, the rights of way cost the Chicago route twice as much in proportion.

\begin{remark}[Round trips and halves]\label{rem:nw:long-haul-fibre:half}
Halving a published round trip assumes both directions take the same path and equipment. On a cable pair they usually do; on a network whose
two directions ride different fibres or routes they need not, and chapter~4's asymmetry returns: a round trip cannot tell the two
directions apart.
\end{remark}

\section{Dark fibre against leased wavelengths}

\begin{definition}[Dark fibre, lit service, wavelength service]\label{def:nw:long-haul-fibre:dark}
\emph{Dark fibre}\index{dark fibre} is fibre sold or leased without equipment: the buyer lights it with its own transponders and amplifiers and
controls everything that affects latency. A \emph{lit service}\index{lit service} is capacity sold on the seller's equipment: a circuit of a
given speed between two points. A \emph{wavelength service}\index{wavelength service} is a lit service that dedicates one wavelength of the
seller's fibre system to the buyer.
\end{definition}

\begin{definition}[Indefeasible right of use]\label{def:nw:long-haul-fibre:iru}
An \emph{indefeasible right of use}\index{indefeasible right of use} (IRU) is a long-term contract giving its holder the exclusive right to use
specified fibres or capacity of a network for a fixed term, usually paid up front, without transferring ownership of the cable, and usually
with a share of maintenance charged to the holder.
\end{definition}

An internet carrier's annual report shows the form in practice: most of its \omterm{def:nw:long-haul-fibre:dark}{dark fibre} is held ``in the form of long-term leases under
indefeasible rights of use'', it relies on the fibre's owner to maintain it, and it amortises the IRUs over generally 15 to 20 years. For a
trading firm, the choice between \omterm{def:nw:long-haul-fibre:dark}{dark fibre} and a \omterm{def:nw:long-haul-fibre:dark}{lit service} is a choice between control and cost. \omterm{def:nw:long-haul-fibre:dark}{Dark fibre} lets the firm choose its own
equipment (light FEC, electronic \omterm{def:nw:long-haul-fibre:disp}{dispersion compensation}, no switches in the path) and costs an IRU, the equipment and yearly maintenance; a
\omterm{def:nw:long-haul-fibre:dark}{lit service} costs a monthly fee and gives the seller's latency.

\begin{omfigure}
\begin{tikzpicture}
\begin{axis}[width=11cm,height=5.8cm,xmin=0,xmax=20,ymin=0,ymax=23,xlabel={years},ylabel={cumulative cost (million USD)},
  tick label style={font=\footnotesize},label style={font=\small},grid=major,grid style={gray!25},legend pos=north west,
  legend style={font=\footnotesize,cells={anchor=west}},table/col sep=comma]
\addplot[omDef,thick,no markers] table[x=year,y=dark_m]{figdata/networks/13-long-haul-fibre/costs.csv};
\addplot[omThm,thick,dashed,no markers] table[x=year,y=lit_m]{figdata/networks/13-long-haul-fibre/costs.csv};
\legend{dark fibre (IRU and equipment up front; O\&M yearly),lit wavelength (monthly fee)}
\end{axis}
\end{tikzpicture}

\omcaption{Illustrative prices, not quotes: the cumulative cost of a long-haul route bought as \omterm{def:nw:long-haul-fibre:dark}{dark fibre} (a 6-million IRU and 1.5~million of
equipment up front, 0.3~million a year of operations and maintenance) and leased as a wavelength (90\,000 a month). The lines cross after 9.6
years. Data: \texttt{fig\_fibre.py}, \texttt{nw\_fibre.costs()}.}\label{fig:nw:long-haul-fibre:costs}
\end{omfigure}

The break-even arithmetic is simple and its inputs are not: IRU prices and lit fees are negotiated, not published, which is why the figure uses
illustrative numbers. With them, \omterm{def:nw:long-haul-fibre:dark}{dark fibre} costs 0.675~million a year spread over twenty years against 1.08~million for the lease, and pays
back its up-front outlay after 9.6 years. A firm that expects to keep the route for less than that, or whose strategy may move to microwave,
leases; a firm that needs to control every microsecond of the path buys, whatever the break-even says.

\omcode{code/firm/fibreroute/firm_fibreroute.py}{82}{96}{The dark-fibre and lit-service cost model: annual costs and the break-even term of buying
against leasing.}\label{lst:nw:long-haul-fibre:costs}

\section{Faster glass: hollow-core fibre}

\begin{definition}[Hollow-core fibre]\label{def:nw:long-haul-fibre:hcf}
A \emph{hollow-core fibre}\index{hollow-core fibre} guides light in a core of air (or vacuum), held by a microstructure of thin glass membranes
around it, so that light travels at nearly $c_0$ instead of at $c_0/1.46$.
\end{definition}

For decades \omterm{def:nw:long-haul-fibre:hcf}{hollow-core fibre} was faster and useless: it lost too much light to cover long distances. A 2025 paper from Microsoft's fibre team
and the University of Southampton reports a \omterm{def:nw:long-haul-fibre:hcf}{hollow-core fibre} with a loss of \qty{0.091}{\decibel} per kilometre at \qty{1550}{\nano\metre},
below the \qty{0.14}{\decibel} of the best solid silica fibre, with transmission speeds 50\% higher, ``a 30\% reduced latency'' and six times
less \omterm{def:nw:long-haul-fibre:disp}{chromatic dispersion} than standard telecom fibre. If the loss holds in cabled routes, the fibre floor moves to nearly the air floor.

\Cref{fig:nw:long-haul-fibre:components} rebuilds the 2010 route four ways with the model (a \omterm{def:nw:the-north-american-map:floor}{group index} of 1.003 for the hollow core is an
assumption consistent with the paper's 30\%). As inverted it takes \qty{6.650}{\milli\second}. With a dispersion-compensating spool of a fifth of
each span, it would take \qty{7.976}{\milli\second}: optical compensation costs more than a millisecond on this route, every reason for a low-latency
route to compensate electronically. In \omterm{def:nw:long-haul-fibre:hcf}{hollow-core fibre} on the same path it would take \qty{4.569}{\milli\second}, \qty{2.08}{\milli\second}
less; in standard fibre along the geodesic, \qty{5.528}{\milli\second}. The new glass would save more than any straightening could.

\begin{omfigure}
\begin{tikzpicture}
\begin{axis}[xbar stacked,width=11.5cm,height=5.4cm,bar width=10pt,xmin=0,xmax=8.6,ytick=data,
  yticklabels from table={figdata/networks/13-long-haul-fibre/components.csv}{variant},y dir=reverse,
  xlabel={one-way latency, 350 Cermak to Carteret (ms)},tick label style={font=\footnotesize},label style={font=\small},
  grid=major,grid style={gray!25},legend columns=4,legend style={font=\footnotesize,at={(0.5,-0.3)},anchor=north,draw=none,
  /tikz/every even column/.append style={column sep=6pt}},table/col sep=comma]
\addplot[fill=omDef!60,draw=omDef] table[y expr=\coordindex,x=glass]{figdata/networks/13-long-haul-fibre/components.csv};
\addplot[fill=omDef!25,draw=omDef] table[y expr=\coordindex,x=path]{figdata/networks/13-long-haul-fibre/components.csv};
\addplot[fill=omMeth!40,draw=omMeth] table[y expr=\coordindex,x=dcf]{figdata/networks/13-long-haul-fibre/components.csv};
\addplot[fill=omThm!60,draw=omThm] table[y expr=\coordindex,x=equipment]{figdata/networks/13-long-haul-fibre/components.csv};
\legend{glass along the geodesic,extra path,DCF,equipment}
\end{axis}
\end{tikzpicture}

\omcaption{The 2010 Chicago--New Jersey route rebuilt four ways by \texttt{firm.fibreroute}: as inverted from its quoted latency, with optical
\omterm{def:nw:long-haul-fibre:disp}{dispersion compensation}, in \omterm{def:nw:long-haul-fibre:hcf}{hollow-core fibre} on the same path, and in standard fibre along the geodesic. Equipment times, the DCF length and the
hollow-core index are the model's assumptions. Data: \texttt{fig\_fibre.py}, \texttt{nw\_fibre.variants()}.}\label{fig:nw:long-haul-fibre:components}
\end{omfigure}

\begin{method}[Choosing a long-haul fibre route]\label{met:nw:long-haul-fibre:choose}
\begin{enumerate}
\item Compute the floor between the two buildings (chapter~10) and ask every seller for a one-way, rack-to-rack latency, with how it was
  measured.
\item Invert each figure into a path factor with a stated equipment assumption; a factor far above the others means a long path or heavy
  equipment, and the seller should say which.
\item Ask how dispersion is compensated and what FEC is used; ask for route diversity (chapter~15) if the route is to carry more than one
  strategy's traffic.
\item Price \omterm{def:nw:long-haul-fibre:dark}{dark fibre} against a \omterm{def:nw:long-haul-fibre:dark}{lit service} over the expected life of the strategy, not of the fibre.
\item Measure the route before and after you depend on it (chapters~4 and~5).
\end{enumerate}
\end{method}

\section{Tutorial: taking a route apart}\label{tut:nw:long-haul-fibre:tut}

\begin{tutorial}
\textbf{Goal.} Model a long-haul route as spans, invert published latencies, and see what \omterm{def:nw:long-haul-fibre:disp}{dispersion compensation} and \omterm{def:nw:long-haul-fibre:hcf}{hollow-core fibre} change.
\textbf{End state:} \cref{fig:nw:long-haul-fibre:sensitivity,fig:nw:long-haul-fibre:costs,fig:nw:long-haul-fibre:components} and
\cref{tab:nw:long-haul-fibre:published}.
\begin{tutsteps}
\item \textbf{The published figures.} \texttt{nw\_fibre.PUBLISHED} holds each route's end points and its one-way figure, from the ledger;
  \texttt{published()} inverts them (\cref{lst:nw:long-haul-fibre:components}) with the equipment of \texttt{EQUIP}.
\item \textbf{Sensitivity.} \texttt{sensitivity()} repeats the inversion for equipment times from 0 to \qty{200}{\micro\second}.
\item \textbf{Variants.} \texttt{variants()} rebuilds the route with \texttt{firm\_fibreroute.build} and \texttt{with\_index}.
\item \textbf{Costs.} \texttt{costs()} and \texttt{break\_even\_years} (\cref{lst:nw:long-haul-fibre:costs}) compare buying and leasing with
  the illustrative \texttt{PRICES}.
\end{tutsteps}
\textbf{What to change next.} Replace the equipment assumptions with a seller's documented figures; add a terrestrial tail to each end of the
Atlantic route and see how much of its excess the tails explain.
\end{tutorial}

\section{Build: the fibre route model}\label{bld:nw:long-haul-fibre:fibreroute}

\begin{build}
\bfield{Purpose} Long-haul fibre routes as data and arithmetic: their latency by component, what a published figure implies, and what buying or
leasing them costs; used by chapters~14 (the fibre back-up of a radio route), 15 (buying) and 29 (the plan).
\bfield{Interface} \texttt{firm\_fibreroute}: \texttt{Span(km, n\_g, amp\_us, dcf\_km)}, \texttt{Route(name, geodesic\_km, spans, terminal\_us)},
\texttt{build}, \texttt{components}, \texttt{invert}, \texttt{with\_index}, \texttt{dark}, \texttt{lit}, \texttt{break\_even\_years}.
\bfield{Rules} A path is never shorter than its geodesic; every equipment time is a stated parameter; published latencies are one-way and say
how they were derived.
\bfield{Acceptance tests} \texttt{code/firm/fibreroute/tests/}: components by hand, inversion round trip, the hollow-core ratio, a DCF spool, and
the break-even arithmetic, including a lease that is always cheaper.
\bfield{Stretch} Routes as polylines over \texttt{firm.geomap} coordinates; loss and amplifier placement from attenuation.
\end{build}

\begin{omsources}
\item G.~Laughlin, A.~Aguirre and J.~Grundfest, ``Information transmission between financial markets in Chicago and New York'', \emph{Financial
Review} 49(2) (2014).
\item M.~Petrovich, E.~Numkam Fokoua et al., ``Broadband optical fibre with an attenuation lower than 0.1 decibel per kilometre'', \emph{Nature
Photonics} (2025).
\item Corning, SMF-28 ULL product information (2020); Equinix, release on Hibernia Express (3 November 2015); Cogent Communications, Form 10-K
for 2022.
\end{omsources}

\section{Exercises}

\begin{exercise}[$\star$]\label{exo:nw:long-haul-fibre:1}
At \qty{0.16}{\decibel} per kilometre, what fraction of the light is left after an 80-kilometre span?
\end{exercise}

\begin{exercise}[$\star$]\label{exo:nw:long-haul-fibre:2}
Using \cref{prop:nw:long-haul-fibre:invert}, how much glass does a route published at \qty{6.55}{\milli\second} contain if its equipment takes
\qty{21.7}{\micro\second}?
\end{exercise}

\begin{exercise}[$\star$]\label{exo:nw:long-haul-fibre:3}
What is the difference between \omterm{def:nw:long-haul-fibre:dark}{dark fibre} and a \omterm{def:nw:long-haul-fibre:dark}{wavelength service}, from the point of view of latency?
\end{exercise}

\begin{exercise}[$\star\star$]\label{exo:nw:long-haul-fibre:4}
Why does optical \omterm{def:nw:long-haul-fibre:disp}{dispersion compensation} cost latency, and electronic compensation not (or much less)?
\end{exercise}

\begin{exercise}[$\star\star$]\label{exo:nw:long-haul-fibre:5}
The Atlantic route's path factor is 1.09 and the Chicago route's 1.20. Give two reasons.
\end{exercise}

\begin{exercise}[$\star\star$]\label{exo:nw:long-haul-fibre:6}
A seller quotes a round trip. What must you ask before halving it?
\end{exercise}

\begin{exercise}[$\star\star\star$]\label{exo:nw:long-haul-fibre:7}
\emph{Coding.} With \texttt{firm.fibreroute}, find the path factor at which a hollow-core route between Cermak and Carteret would match a
microwave route at a \omterm{def:nw:the-north-american-map:factor}{route factor} of 1.011 in air.
\end{exercise}

\begin{exercise}[$\star\star\star$]\label{exo:nw:long-haul-fibre:8}
\emph{Find the flaw.} ``Our new fibre route is only 3\% longer than the geodesic, so it will beat the microwave networks, whose towers zig-zag.''
\end{exercise}

\section{Problem: Thirteen Milliseconds}

\begin{problem}[{Weekend problem --- taking the 2010 Chicago--New Jersey route apart}]\label{pb:nw:long-haul-fibre:1}
The 2010 fibre route from 350 Cermak to Carteret was quoted at \qty{6.65}{\milli\second} one way. Use the model's assumptions: \qty{10}{\micro\second}
per terminal, an amplifier every 80~kilometres at \qty{0.1}{\micro\second}, standard fibre of \omterm{def:nw:the-north-american-map:floor}{group index} 1.462.

\medskip\noindent\textbf{Part I --- The floor.}
\begin{enumerate}
\item What are the \omterm{def:nw:the-north-american-map:geodesic}{geodesic distance} and the floors in vacuum and in fibre between the two buildings?
\item What is the round trip quoted, and the round-trip floor in fibre?
\item How much of the one-way figure is above the fibre floor?
\item Why is the fibre floor, not the vacuum floor, the right comparison for a fibre route?
\end{enumerate}

\medskip\noindent\textbf{Part II --- The split.}
\begin{enumerate}[resume]
\item How much glass does the route contain, and how many spans?
\item What is the equipment time in the model, and how much of the excess is equipment?
\item How much is extra path, in microseconds and in kilometres?
\item What happens to the split if the equipment takes \qty{200}{\micro\second}?
\end{enumerate}

\medskip\noindent\textbf{Part III --- Other glass.}
\begin{enumerate}[resume]
\item What would optical \omterm{def:nw:long-haul-fibre:disp}{dispersion compensation} of a fifth of each span add?
\item What would the same path take in \omterm{def:nw:long-haul-fibre:hcf}{hollow-core fibre} of \omterm{def:nw:the-north-american-map:floor}{group index} 1.003?
\item And standard fibre along the geodesic?
\item How do these compare with the 2016 microwave route of chapter~10 (\qty{3.982}{\milli\second} from Aurora)?
\end{enumerate}

\medskip\noindent\textbf{Part IV --- The verdict.}
\begin{enumerate}[resume]
\item State the \emph{named result}: the split of the 2010 round trip into glass along the geodesic, extra path and equipment, and the milliseconds
  a \omterm{def:nw:long-haul-fibre:hcf}{hollow-core fibre} on the same path would save.
\item Which part of the split depends most on the model's assumptions?
\item Why did the route lose its lead within two years?
\item What did it keep?
\item Would you buy it as \omterm{def:nw:long-haul-fibre:dark}{dark fibre} or lease a wavelength on it?
\item What would a hollow-core route need to beat the microwave networks?
\item What does the Atlantic route's path factor say about where the Chicago route lost its time?
\item In one sentence: where do the thirteen milliseconds go?
\end{enumerate}
\end{problem}

\section{Interview questions}

\begin{interviewq}[$\star$ \iqroles{developer}]\label{iq:nw:long-haul-fibre:1}
How fast does light travel in fibre, and what is the latency per 1\,000 kilometres?
\end{interviewq}

\begin{interviewq}[$\star\star$ \iqroles{developer}]\label{iq:nw:long-haul-fibre:2}
What adds latency to a long-haul fibre route besides the length of the glass?
\end{interviewq}

\begin{interviewq}[$\star\star$ \iqroles{developer, trader}]\label{iq:nw:long-haul-fibre:3}
Why did microwave beat fibre between Chicago and New Jersey, and why is fibre still used there?
\end{interviewq}

\begin{interviewq}[$\star\star$ \iqroles{developer}]\label{iq:nw:long-haul-fibre:4}
What is an IRU, and when would a trading firm buy \omterm{def:nw:long-haul-fibre:dark}{dark fibre} rather than lease a wavelength?
\end{interviewq}

\begin{interviewq}[$\star\star$ \iqroles{developer, researcher}]\label{iq:nw:long-haul-fibre:5}
What is \omterm{def:nw:long-haul-fibre:hcf}{hollow-core fibre}, and what would it change for latency-sensitive trading?
\end{interviewq}

\begin{interviewq}[$\star\star\star$ \iqroles{developer}]\label{iq:nw:long-haul-fibre:6}
A seller claims a fibre route between two cities that is faster than your computed fibre floor. What could explain it?
\end{interviewq}
